% ============================================================================
%  Syllabus — Qualité et tests logiciels (L2)
%  Dr. Yvan GUIFO FODJO — EFREI Paris — 2025–2026
%  Compilation : pdflatex (Overleaf)
% ============================================================================
\documentclass[11pt,a4paper]{article}

% ---- Encodage et langue ----------------------------------------------------
\usepackage[utf8]{inputenc}
\usepackage[T1]{fontenc}
\usepackage[french]{babel}
\usepackage{lmodern}

% ---- Mise en page ----------------------------------------------------------
\usepackage[a4paper,margin=2.3cm,top=2.4cm,bottom=2.4cm]{geometry}
\setlength{\parindent}{0pt}
\setlength{\parskip}{0.45em}
\setlength{\emergencystretch}{2em}

% ---- Couleurs (alignées avec SQ_Seance1_v2.tex) ----------------------------
\usepackage{xcolor}
\definecolor{accent}{HTML}{3B82F6}
\definecolor{accent2}{HTML}{0EA5A4}
\definecolor{bgsoft}{HTML}{F1F5F9}
\definecolor{linecolor}{HTML}{CBD5E1}
\definecolor{textsoft}{HTML}{475569}
\definecolor{bgaccent}{HTML}{E0EBFB}
\definecolor{bgaccent2}{HTML}{D8F0EE}

% ---- Titres ----------------------------------------------------------------
\usepackage{titlesec}
\titleformat{\section}
  {\color{accent}\Large\bfseries}
  {\thesection.}{0.6em}{}[\vspace{2pt}{\color{accent}\titlerule[1.2pt]}]
\titleformat{\subsection}
  {\color{accent2}\large\bfseries}
  {\thesubsection}{0.5em}{}
\titlespacing*{\section}{0pt}{1.4em}{0.6em}
\titlespacing*{\subsection}{0pt}{1.0em}{0.3em}

% ---- Tableaux --------------------------------------------------------------
\usepackage{booktabs}
\usepackage{array}
\usepackage{tabularx}
\usepackage{colortbl}
\renewcommand{\arraystretch}{1.25}
\arrayrulecolor{linecolor}
\newcolumntype{L}[1]{>{\raggedright\arraybackslash}p{#1}}
\newcolumntype{C}[1]{>{\centering\arraybackslash}p{#1}}

% ---- Listes ----------------------------------------------------------------
\usepackage{enumitem}
\setlist{itemsep=2pt,topsep=4pt,parsep=0pt,leftmargin=1.4em}

% ---- Hyperliens et URLs ---------------------------------------------------
\PassOptionsToPackage{hyphens}{url}
\usepackage{hyperref}
\usepackage{xurl}
\Urlmuskip=0mu plus 1mu
\hypersetup{
    colorlinks=true,
    linkcolor=accent,
    urlcolor=accent2,
    citecolor=accent,
    pdftitle={Syllabus - Qualite et tests logiciels - L2},
    pdfauthor={Dr. Yvan GUIFO FODJO}
}

% ---- En-tête / pied de page -----------------------------------------------
\usepackage{fancyhdr}
\pagestyle{fancy}
\fancyhf{}
\renewcommand{\headrulewidth}{0pt}
\renewcommand{\footrulewidth}{0.4pt}
\fancyfoot[L]{\color{textsoft}\footnotesize Dr. Yvan GUIFO FODJO --- EFREI Paris --- 2025--2026}
\fancyfoot[R]{\color{textsoft}\footnotesize Page \thepage}

% ---- Encadrés --------------------------------------------------------------
\usepackage[most]{tcolorbox}
\newtcolorbox{noteinfo}{
    enhanced,
    colback=bgaccent,
    colframe=accent,
    boxrule=0pt,
    leftrule=3pt,
    arc=2pt,
    left=8pt, right=8pt, top=6pt, bottom=6pt,
    fontupper=\small
}
\newtcolorbox{noteeval}{
    enhanced,
    colback=bgaccent2,
    colframe=accent2,
    boxrule=0pt,
    leftrule=3pt,
    arc=2pt,
    left=8pt, right=8pt, top=6pt, bottom=6pt,
    fontupper=\small
}

% ============================================================================
\begin{document}

% ---- BLOC TITRE -------------------------------------------------------------
\begin{center}
{\color{accent}\rule{0.22\linewidth}{1.4pt}}\\[0.4em]
{\Huge\bfseries Qualité et tests logiciels}\\[0.4em]
{\large\color{textsoft}Syllabus de cours --- Niveau L2 ingénierie informatique}\\[0.8em]
\begin{tcolorbox}[
    enhanced, colback=bgsoft, colframe=linecolor, boxrule=0.6pt,
    arc=3pt, left=10pt, right=10pt, top=8pt, bottom=8pt,
    width=0.86\linewidth
]
\centering
{\large\bfseries Dr. Yvan GUIFO FODJO}\\[0.2em]
{\color{textsoft} EFREI Paris \quad|\quad Format CTP \quad|\quad 15h \quad|\quad 3 séances de 5h}\\[0.2em]
{\small\color{textsoft} Année universitaire 2025--2026 \quad|\quad \texttt{yvan.guifo-fodjo@efrei.fr}}
\end{tcolorbox}
\end{center}
\vspace{0.4em}

% ============================================================================
\section{Identification du cours}

\begin{tabularx}{\linewidth}{@{}>{\bfseries\color{accent}}p{4.5cm} X@{}}
\toprule
Intitulé & Qualité et tests logiciels \\
Niveau & L2 --- Ingénierie informatique \\
Volume horaire & 15 h en présentiel --- 3 séances de 5 h \\
Format & CTP (Cours-Travaux Pratiques intégrés) \\
Charge personnelle estimée & 8 à 10 h (préparation, révisions QCM) \\
Année universitaire & 2025--2026 \\
Enseignant & Dr. Yvan GUIFO FODJO --- EFREI Paris \\
Contact & \texttt{yvan.guifo-fodjo@efrei.fr} \\
\bottomrule
\end{tabularx}

% ============================================================================
\section{Prérequis}

\begin{itemize}
    \item \textbf{Programmation Python} : syntaxe, fonctions, classes, exceptions, modules.
    \item \textbf{Programmation Java} : syntaxe, classes, héritage, exceptions, Maven ou Gradle (notions).
    \item \textbf{Git} : \texttt{clone}, \texttt{add}, \texttt{commit}, \texttt{push} (manipulation de base).
    \item \textbf{Ligne de commande} : navigation, exécution de scripts.
\end{itemize}

Aucun prérequis en théorie des tests n'est attendu.

% ============================================================================
\section{Description et raison d'être}

Tout logiciel non trivial contient des défauts. La capacité à \textbf{prévenir, détecter et corriger} ces défauts distingue un développeur débutant d'un ingénieur logiciel. Ce cours introduit les \textbf{fondements théoriques} et les \textbf{pratiques professionnelles} du test logiciel, en s'appuyant sur deux écosystèmes industriels majeurs : \textbf{Python (pytest)} et \textbf{Java (JUnit~5)}.

Le format CTP entrelace systématiquement exposé conceptuel et manipulation : chaque notion introduite est immédiatement mise en \oe uvre sur un \emph{fil rouge} --- une mini-application de \textbf{gestion de librairie en ligne} --- qui s'enrichit séance après séance.

% ============================================================================
\section{Objectifs d'apprentissage}

\textit{Formulés selon la taxonomie de Bloom révisée (Anderson \& Krathwohl, 2001).}

À l'issue du cours, l'étudiant·e sera capable de :

\begin{tabularx}{\linewidth}{@{}C{0.7cm} X C{2.6cm}@{}}
\toprule
\textbf{\#} & \textbf{Objectif} & \textbf{Niveau Bloom} \\
\midrule
O1 & \textbf{Définir} la terminologie ISTQB (erreur, défaut, panne, V\&V) et \textbf{distinguer} vérification et validation & Comprendre (2) \\
O2 & \textbf{Identifier} les caractéristiques de qualité ISO/IEC 25010 affectées dans un cas concret & Comprendre (2) \\
O3 & \textbf{Appliquer} le patron AAA pour écrire des tests unitaires en \texttt{pytest} et \texttt{JUnit~5} & Appliquer (3) \\
O4 & \textbf{Concevoir} des cas de test par partitions d'équivalence et valeurs limites & Appliquer (3) / Analyser (4) \\
O5 & \textbf{Mesurer} et \textbf{interpréter} la couverture de code avec \texttt{coverage.py} et \texttt{JaCoCo} & Appliquer (3) \\
O6 & \textbf{Utiliser} des doubles de test (mocks, stubs, fakes) avec \texttt{unittest.mock} et \texttt{Mockito} & Appliquer (3) \\
O7 & \textbf{Pratiquer} le cycle TDD \emph{red-green-refactor} sur un kata court & Appliquer (3) \\
O8 & \textbf{Mettre en place} un pipeline d'intégration continue minimal avec GitHub Actions & Appliquer (3) \\
O9 & \textbf{Critiquer} une suite de tests existante au regard des principes FIRST & Analyser (4) \\
\bottomrule
\end{tabularx}

% ============================================================================
\section{Alignement constructif (Biggs, 1996)}

\begin{tabularx}{\linewidth}{@{}C{1.3cm} X L{4.2cm}@{}}
\toprule
\textbf{Objectif} & \textbf{Activité CTP} & \textbf{Évaluation} \\
\midrule
O1, O2 & Exercices d'analyse de cas (S1) & QCM final \\
O3 & TP guidés pytest / JUnit (S1) & TP noté S2 + QCM final \\
O4 & Atelier conception de tests (S2) & TP noté S2 + QCM final \\
O5 & TP couverture (S2) & TP noté S2 + QCM final \\
O6 & TP mocking (S2) & QCM final \\
O7 & Kata TDD (S3) & QCM final \\
O8 & TP GitHub Actions (S3) & QCM final \\
O9 & Revue de tests entre pairs (S3) & QCM final \\
\bottomrule
\end{tabularx}

% ============================================================================
\section{Programme détaillé}

% ---- SÉANCE 1 --------------------------------------------------------------
\subsection{Séance 1 (5 h) --- Fondements et premiers tests}

\begin{tabularx}{\linewidth}{@{}C{1.0cm} C{1.6cm} X@{}}
\toprule
\textbf{Bloc} & \textbf{Durée} & \textbf{Contenu} \\
\midrule
1.1 & 40 min & Qualité logicielle, ISO/IEC 25010, coût des défauts \\
1.2 & 30 min & Vocabulaire ISTQB : erreur / défaut / panne, V\&V \\
1.3 & 50 min & Pyramide de tests (Cohn 2009), principes FIRST \\
1.4 & 60 min & TP1 --- Premiers tests \texttt{pytest}, patron AAA \\
1.5 & 45 min & Cas limites, \texttt{@pytest.mark.parametrize}, \texttt{pytest.approx} \\
1.6 & 60 min & TP2 --- Tests \texttt{JUnit~5}, \texttt{@BeforeEach}, fixtures \\
Pauses & 15 min & --- \\
\bottomrule
\end{tabularx}

\begin{noteinfo}
\textbf{Livrable étudiant (non noté)} : suite de tests pour le module \texttt{LibraireService} (Python + Java). Sert de préparation au TP noté de S2.
\end{noteinfo}

% ---- SÉANCE 2 --------------------------------------------------------------
\subsection{Séance 2 (5 h) --- Conception de tests, couverture et TP noté}

\begin{tabularx}{\linewidth}{@{}C{1.0cm} C{1.6cm} X@{}}
\toprule
\textbf{Bloc} & \textbf{Durée} & \textbf{Contenu} \\
\midrule
2.1 & 30 min & Boîte noire vs boîte blanche ; rappel de S1 \\
2.2 & 45 min & Partitions d'équivalence et valeurs limites (Beizer 1990) \\
2.3 & 40 min & Tables de décision et tests basés sur états \\
2.4 & 50 min & TP3 (guidé) --- Conception de tests sur le module \texttt{Panier} \\
2.5 & 40 min & Couverture de code : instruction, branche, condition (Ammann \& Offutt 2017) \\
2.6 & 40 min & TP4 (guidé) --- \texttt{coverage.py} et \texttt{JaCoCo}, lecture de rapports HTML \\
2.7 & 35 min & Test doubles : \emph{dummies, stubs, fakes, mocks, spies} (Meszaros 2007) \\
\rowcolor{bgaccent2!50}
\textbf{2.8} & \textbf{60 min} & \textbf{TP noté individuel sur machine} \\
Pauses & 20 min & --- \\
\bottomrule
\end{tabularx}

\begin{noteeval}
\textbf{TP noté (S2, bloc 2.8 --- 60 min)} : sujet inédit, code Python fourni. L'étudiant écrit une suite de tests respectant AAA + parametrize, atteint un seuil de couverture imposé et identifie au moins un défaut introduit volontairement. Grille analytique remise en début de séance.
\end{noteeval}

% ---- SÉANCE 3 --------------------------------------------------------------
\subsection{Séance 3 (5 h) --- TDD, intégration, CI et examen final}

\begin{tabularx}{\linewidth}{@{}C{1.0cm} C{1.6cm} X@{}}
\toprule
\textbf{Bloc} & \textbf{Durée} & \textbf{Contenu} \\
\midrule
3.1 & 25 min & Test-Driven Development : cycle \emph{red-green-refactor} (Beck 2002) \\
3.2 & 60 min & TP5 --- Kata FizzBuzz en TDD (étape par étape, en binôme) \\
3.3 & 25 min & Refactoring guidé par les tests \\
3.4 & 35 min & Tests d'intégration : base de données en mémoire (SQLite, H2) \\
3.5 & 50 min & TP6 --- Intégration continue avec GitHub Actions : workflow \texttt{.yml}, badges \\
3.6 & 20 min & Mention : analyse statique (\texttt{pylint}, SpotBugs), tests de mutation, fuzzing \\
3.7 & 15 min & Synthèse, conseils de révision pour le QCM \\
\rowcolor{bgaccent2!50}
\textbf{3.8} & \textbf{90 min} & \textbf{Examen final QCM individuel} (sans documents) \\
Pauses & 20 min & --- \\
\bottomrule
\end{tabularx}

\begin{noteeval}
\textbf{Examen final (S3, bloc 3.8 --- 1h30)} : QCM de 30 à 40 questions couvrant l'ensemble du programme, avec distracteurs plausibles. Pas de pénalité pour réponse fausse (à confirmer selon politique EFREI).
\end{noteeval}

% ============================================================================
\section{Méthodes pédagogiques}

\begin{itemize}
    \item \textbf{Format CTP} : alternance d'exposés courts (15--25 min) et de manipulations guidées (40--60 min), conformément aux résultats de Freeman et al.\ (2014, PNAS, DOI : \href{https://doi.org/10.1073/pnas.1319030111}{10.1073/pnas.1319030111}) sur l'efficacité de l'apprentissage actif.
    \item \textbf{Fil rouge} : application \emph{Librairie en ligne} (catalogue, panier, paiement simulé) enrichie sur les 3 séances.
    \item \textbf{Pair programming} ponctuel en TP pour favoriser l'entraide.
    \item \textbf{Feedback immédiat} : correction collective en fin de chaque TP.
\end{itemize}

% ============================================================================
\section{Outils et environnement}

\textit{Versions cibles données à titre indicatif ; alternatives libres signalées en séance.}

\begin{tabularx}{\linewidth}{@{}L{3.4cm} L{2.2cm} X@{}}
\toprule
\textbf{Outil} & \textbf{Version} & \textbf{Usage} \\
\midrule
Python & 3.12+ & Code et tests \\
\texttt{pytest} & 8.x & Framework de test Python \\
\texttt{pytest-cov} & 5.x & Couverture côté Python \\
\texttt{coverage.py} & 7.x & Mesure de couverture \\
Java & 21 (LTS) & Code et tests \\
JUnit & 5.10+ & Framework de test Java \\
Mockito & 5.x & Mocks Java \\
JaCoCo & 0.8.x & Couverture côté Java \\
Maven & 3.9+ & Build Java \\
Git & 2.40+ & Versionnement \\
GitHub Actions & --- & Intégration continue \\
IDE & VS Code ou IntelliJ Community & Au choix de l'étudiant \\
\bottomrule
\end{tabularx}

% ============================================================================
\section{Évaluation}

\begin{tabularx}{\linewidth}{@{}L{4.5cm} C{1.8cm} X C{1.5cm}@{}}
\toprule
\textbf{Composante} & \textbf{Poids} & \textbf{Modalité} & \textbf{Durée} \\
\midrule
TP noté individuel --- S2 & \textbf{40\,\%} & Sur machine, sujet écrit, code Python à tester & 60 min \\
Examen final QCM --- S3 & \textbf{60\,\%} & Sur table ou plateforme, sans documents & 90 min \\
\bottomrule
\end{tabularx}

\subsection{Grille du TP noté de S2 (sur 20)}

\begin{tabularx}{\linewidth}{@{}X C{1.5cm}@{}}
\toprule
\textbf{Critère} & \textbf{Points} \\
\midrule
Application stricte du patron AAA --- nommage explicite & 4 \\
Couverture des cas nominaux et des cas limites & 5 \\
Usage pertinent de \texttt{parametrize} ou de tests équivalents & 3 \\
Détection du ou des défaut(s) introduit(s) & 4 \\
Couverture de code $\geq$ seuil imposé & 3 \\
Lisibilité, indépendance et rapidité des tests (FIRST) & 1 \\
\midrule
\textbf{Total} & \textbf{20} \\
\bottomrule
\end{tabularx}

\subsection{Composition indicative du QCM final}

30 à 40 questions, équilibrage Bloom :
\begin{itemize}
    \item Niveau Comprendre (O1, O2) : 30\,\% des questions
    \item Niveau Appliquer (O3 à O8) : 50\,\% des questions
    \item Niveau Analyser (O9 + lecture de code) : 20\,\% des questions
\end{itemize}

% ============================================================================
\section{Politique d'intégrité académique}

\begin{itemize}
    \item Le plagiat de code (entre étudiants ou depuis le web) entraîne la note de 0 et un signalement.
    \item L'utilisation d'assistants (Copilot, Claude, ChatGPT) est \textbf{autorisée pour explorer et apprendre}, \textbf{interdite pendant le TP noté et l'examen}. Tout code rendu doit être compris et défendable à l'oral.
    \item Toute soumission générée par IA hors contexte autorisé doit être explicitement déclarée.
\end{itemize}

% ============================================================================
\section{Bibliographie}

\subsection{Références principales}

\begin{enumerate}
    \item Anderson, L.~W.\ et Krathwohl, D.~R.\ (dir.). \emph{A Taxonomy for Learning, Teaching, and Assessing: A Revision of Bloom's Taxonomy of Educational Objectives}. New York : Longman, 2001. ISBN 978-0801319037.
    \item Ammann, P.\ et Offutt, J. \emph{Introduction to Software Testing}. 2\textsuperscript{e}~éd. Cambridge University Press, 2017. ISBN 978-1107172012.
    \item Beck, K. \emph{Test-Driven Development: By Example}. Addison-Wesley, 2002. ISBN 978-0321146533.
    \item Biggs, J. «~Enhancing teaching through constructive alignment~». \emph{Higher Education}, vol.~32, n\textsuperscript{o}~3, 1996, p.~347--364. DOI : \href{https://doi.org/10.1007/BF00138871}{10.1007/BF00138871}.
    \item Cohn, M. \emph{Succeeding with Agile: Software Development Using Scrum}. Addison-Wesley, 2009. ISBN 978-0321579362 (chapitre~16, \emph{The Test Automation Pyramid}).
    \item Meszaros, G. \emph{xUnit Test Patterns: Refactoring Test Code}. Addison-Wesley, 2007. ISBN 978-0131495050.
\end{enumerate}

\subsection{Standards et référentiels}

\begin{enumerate}\setcounter{enumi}{6}
    \item ISO/IEC 25010:2011. \emph{Systems and software engineering --- Systems and software Quality Requirements and Evaluation (SQuaRE) --- System and software quality models}. Genève : ISO, 2011. \url{https://www.iso.org/standard/35733.html}
    \item ISO/IEC/IEEE 29119 (séries 1 à 5). \emph{Software and systems engineering --- Software testing}. \url{https://www.iso.org/standard/81291.html}
    \item IEEE Computer Society. \emph{SWEBOK v3.0 --- Guide to the Software Engineering Body of Knowledge}, chapitre~4 \emph{Software Testing}. IEEE, 2014. \url{https://www.computer.org/education/bodies-of-knowledge/software-engineering}
    \item ISTQB. \emph{Foundation Level Syllabus v4.0}. International Software Testing Qualifications Board, 2023. \url{https://www.istqb.org/certifications/certified-tester-foundation-level}
\end{enumerate}

\subsection{Référence complémentaire}

\begin{enumerate}\setcounter{enumi}{10}
    \item Freeman, S.\ et al.\ «~Active learning increases student performance in science, engineering, and mathematics~». \emph{Proceedings of the National Academy of Sciences}, vol.~111, n\textsuperscript{o}~23, 2014, p.~8410--8415. DOI : \href{https://doi.org/10.1073/pnas.1319030111}{10.1073/pnas.1319030111}.
\end{enumerate}

\subsection{Documentation technique en ligne}

\begin{enumerate}\setcounter{enumi}{11}
    \item \emph{pytest documentation}. \url{https://docs.pytest.org/}
    \item \emph{JUnit 5 User Guide}. \url{https://junit.org/junit5/docs/current/user-guide/}
    \item \emph{Mockito Reference}. \url{https://site.mockito.org/}
    \item \emph{coverage.py docs}. \url{https://coverage.readthedocs.io/}
    \item \emph{JaCoCo docs}. \url{https://www.jacoco.org/jacoco/trunk/doc/}
\end{enumerate}

\vspace{1em}
\begin{center}
\color{textsoft}\small\itshape
--- Fin du syllabus ---\\
Document susceptible de mises à jour mineures avant la première séance.
\end{center}

\end{document}
